% Copyright (c) 2026 Christian Reitwiessner.
% Authors: Christian Reitwiessner
% Build: /tmp/andreev-pdf-tools/tectonic --keep-logs paper.tex
\documentclass[11pt]{article}
\usepackage[T1]{fontenc}
\usepackage{lmodern,amsmath,amssymb,amsthm,booktabs,microtype,enumitem}
\usepackage[margin=1in]{geometry}
\usepackage[colorlinks=true,allcolors=blue]{hyperref}
\newtheorem{theorem}{Theorem}[section]
\newtheorem{lemma}[theorem]{Lemma}
\newtheorem{proposition}[theorem]{Proposition}
\newtheorem{corollary}[theorem]{Corollary}
\newcommand{\bits}{\{0,1\}}
\newcommand{\ind}[1]{[\,#1\,]}
\newcommand{\THR}{\mathsf{THR}}
\newcommand{\smalltitle}[1]{\paragraph{#1}}
\setlength{\emergencystretch}{1em}
\setlist[itemize]{itemsep=1pt,topsep=3pt}
\setlist[enumerate]{itemsep=1pt,topsep=3pt}
\title{Exact and Approximate Threshold Circuits\\for Andreev's Function}
\author{Christian Reitwiessner}
\date{November 2026}
\begin{document}
\maketitle
\begin{abstract}
\marginpar{\raggedright\scriptsize\textbf{TODO:} The abstract still needs considerable work.}
Andreev's function reads an input table at an address formed by parities
of disjoint input blocks. For the balanced $n$-input family, we construct
exact depth-two threshold circuits with $(5/16+o(1))n^2$ gates, improving
the $O(n^3/\log n)$ upper bound stated by Kane and Williams. Under
uniform inputs, explicit deterministic circuits achieve error at most
$\delta$ with $O(n^{3/2}\sqrt{\log(1/\delta)})$ gates, for
$0<\delta\le1/2$. At fixed $0<\delta<1/2$, the exponent matches
their lower bound up to logarithmic factors. Symmetric truncation of a
Fourier lookup formula makes the approximation errors odd in distinct
table coefficients; a table symmetry then makes these errors orthogonal.
Both constructions use bottom gates whose address weights have one sign
within each block. A restriction-counting argument gives an exact lower
bound of $\Omega(n^2/\log^3 n)$ for this class, with unrestricted real
table and output weights. We also give a size--advantage tradeoff and a
separate exponential bound for full-support, equal-magnitude address weights.
\end{abstract}

\section{Introduction}
A depth-two threshold circuit has bottom threshold gates on the input
bits and one output threshold gate. Each gate outputs one when its
weighted sum meets or exceeds its threshold. We allow arbitrary real
weights and direct input wires to the output. Size counts all gates;
counts of bottom gates add one output gate separately. Wires are edges
with nonzero coefficients after coefficients on each edge are combined.

Fix $k,m\ge1$ and $N=2^k$. The inputs are an $N$-bit table
$x=(x_c)_{c\in\bits^k}$ and $k$ disjoint address blocks $a_i\in\bits^m$.
The ordinary Andreev function reads the table at the address given by
the block parities:
$$
b_i=|a_i|\bmod2,\qquad A_{k,m}(x,a)=x_b,\qquad n=N+km.
$$
Table and address bits are separate inputs. Asymptotic statements use
\emph{floor balance}, $m=\lfloor N/k\rfloor$, so $n\sim2N$.
We also consider \emph{exact balance}, $km=N$, when $k\mid2^k$.
Finite-parameter results state their hypotheses separately.

\paragraph{Results and methods.}
Fourier inversion expresses lookup as a sum of table coefficients times
address parities. Section~\ref{sec:identities} compiles each product directly
into bottom thresholds, giving quadratic exact size; Appendix~\ref{sec:exact}
refines the constant to $5/16$. Section~\ref{sec:approx} truncates the
Fourier terms symmetrically. A table symmetry makes their errors
orthogonal, yielding deterministic approximation circuits. Table~\ref{tab:bounds}
summarizes the main comparisons.

A bottom gate is \emph{sign-coherent} if its weights within each address
block are all nonnegative or all nonpositive; signs may vary by block and
gate. Our upper bounds use such gates. Their exact complexity has exponent
$2$, while their fixed-error complexity has exponent $3/2$, up to logarithms.

\begin{table}[t]
\centering
\small
\setlength{\tabcolsep}{4pt}
\renewcommand{\arraystretch}{1.25}
\begin{tabular}{@{}lll@{}}
\toprule
Setting and resource & Earlier bound~\cite{KW} & Bound here \\
\midrule
Exact, gate upper bound
 & $O(n^3/\log n)$ & $(5/16+o(1))n^2$ \\
Advantage $\varepsilon$, gates
 & $\Omega\!\left(\frac{\varepsilon^3n^{3/2}}{\log^3 n}\right)$
 & $O\!\left(\varepsilon^2n^{3/2}\sqrt{\log(1/\varepsilon)}\right)$ \\
Advantage $\varepsilon$, wires
 & $\Omega\!\left(\frac{\varepsilon^3n^{5/2}}{\log^{7/2}n}\right)$
 & $O\!\left(\varepsilon^2n^{5/2}\sqrt{\log(1/\varepsilon)}\right)$ \\
Exact, coherent lower bound
 & $\Omega(n^{3/2}/\log^3 n)$ & $\Omega(n^2/\log^3 n)$ \\
\bottomrule
\end{tabular}
\caption{Balanced ordinary Andreev lookup. Advantage $\varepsilon$ means
uniform-input agreement at least $1/2+\varepsilon$. The approximation rows
compare earlier lower bounds with our upper bounds, matching up to logarithms
at fixed advantage. Earlier bounds require $\varepsilon\gg\sqrt{\log n/n}$;
our displayed upper bounds hold for $N^{-1/2}\le\varepsilon\le1/8$.
Theorem~\ref{thm:avgupper} also handles larger advantage and gives the error
parameterization. All upper bounds here are sign-coherent. The last row
compares our restricted lower bound with the earlier unrestricted bound.}
\label{tab:bounds}
\end{table}

\paragraph{Relation to earlier work.}
Kane and Williams~\cite[Theorems~1.1--1.2]{KW} give the earlier bounds
in Table~\ref{tab:bounds}. Their Section~5.2 also gives a quadratic bound
in the different $\mathsf{MOD}_3\circ\mathsf{MOD}_2$ basis. Their threshold
construction expands lookup tests into parities and simulates them
separately; ours first groups by address character. Their Section~4
approximates parity by central-window staircases, and Section~6 notes
correlation with table majority. Our approximation adds odd-error
orthogonality and applies it to partial table votes.

Threshold staircases are classical; see Hajnal et al.~\cite[Proposition~2.1]{HMPST}
for symmetric functions and Hansen--Podolskii~\cite[Theorem~5]{HP} for a
related representation using bounded integer forms. Our sign-coherent
lower bound combines an information count under restrictions with the
single-crossing principle used in threshold lower bounds~\cite{IPS}.
It requires no weight bound. The polynomial overhead in simulations
with monotone bottom gates~\cite[Lemma~6]{HP} does not preserve our bounds.
The generalized Andreev function of Chen, Santhanam and
Srinivasan~\cite[Section~2.4]{CSS} uses an encoded table and an extractor
address, unlike the ordinary lookup function here. The more recent
$n^{5/2-\eta}$ gate lower bounds for functions in
$E^{\mathsf{NP}}$~\cite{CTW} also concern other functions; our exponent
gap is specific to ordinary Andreev lookup.

Appendix~\ref{sec:fullsign} gives a separate exponential lower bound
for full-support, equal-magnitude address weights. It uses the classical
parity-order argument of Minsky and Papert~\cite[Section~3.1]{MP}, which
forces a predicate to depend on every parity variable. It does not
improve the unrestricted lower bound. We make no exhaustive priority
claim for the elementary ingredients.

\section{Compiling functions of affine forms}\label{sec:identities}
We first provide a construction that can save circuit layers when
realizing certain functions. Write $[P]$ for the Boolean indicator of a statement $P$.

\begin{lemma}[Affine-form compilation]\label{lem:ladder}
Let $d\in\mathbb N_{\ge1}$, let $L:\bits^d\to\mathbb Z$ be an affine
form, and let $\ell,h\in\mathbb Z$ satisfy $\ell\le L(u)\le h$ for every
$u\in\bits^d$. For every function $g:\{\ell,\ldots,h\}\to\mathbb R$
and every $u\in\bits^d$,
$$
g(L(u))=g(\ell)+\sum_{\theta=\ell+1}^{h}
\bigl(g(\theta)-g(\theta-1)\bigr)[L(u)\ge\theta].
$$
Consequently, a score consisting of an affine term and a sum of such
$g(L)$ terms can be thresholded by a depth-two circuit with direct input
wires to the output. Each nonzero finite difference costs at most one
bottom gate, and identical gates can be shared.
\end{lemma}
\begin{proof}
Fix $u\in\bits^d$. The active comparisons are exactly those with
$\theta\le L(u)$, so the
sum telescopes. Each comparison is a threshold gate on the input bits.
The output threshold absorbs the constants, and direct input wires
supply the affine terms.
\end{proof}

To apply the lemma to lookup, define
$$
X_S=\sum_{c\in\bits^k}(-1)^{S\cdot c+x_c},\qquad
t_S=\sum_{i\in S}|a_i|,\qquad S\subseteq[k],
$$
where $S\cdot c=\sum_{i\in S}c_i$. The table coefficient $X_S$ is an
even integer in $[-N,N]$. It is affine in the input bits because
$(-1)^{x_c}=1-2x_c$, and $t_S$ is affine as well. Character orthogonality
gives the lookup identity
\begin{equation}\label{eq:lookup}
\sum_{S\subseteq[k]}X_S(-1)^{t_S}=N(-1)^{x_b}.
\end{equation}
To verify it, exchange the sums. The coefficient of $(-1)^{x_c}$ becomes
$\prod_{i=1}^k(1+(-1)^{c_i+b_i})$, which is $N$ if $c=b$ and zero
otherwise. In matrix notation, let $v_c=(-1)^{x_c}$ and
$H_{S,c}=(-1)^{S\cdot c}$. Then $X=Hv$, $H^{\mathsf T}H=NI$, and
$v=N^{-1}H^{\mathsf T}X$. Thus $X_S/N$ are the coordinates of the signed
truth table in the parity basis, and \eqref{eq:lookup} recovers its $b$th
entry by Fourier inversion; see also \cite[Chapter~1]{OD}.

\begin{theorem}[A quadratic exact circuit]\label{thm:simple}
For all $k,m\ge1$, $A_{k,m}$ has a sign-coherent circuit with at most
$2N(N-1)+kmN$ bottom gates and one output gate, using polynomially
bounded integer weights. In the balanced family this is
$(3/4+o(1))n^2$ gates.
\end{theorem}
\begin{proof}
To separate the table coefficient from the address parity, put
$\varphi(u)=(-1)^u u$ for integers $u$. If $X$ is even, then
$$X(-1)^t=\varphi(X+t)-\varphi(t).$$
Each nonempty summand of \eqref{eq:lookup} is therefore a difference
of functions of affine forms, with range lengths $2N+m|S|$ and $m|S|$.
The empty-set summand $X_\varnothing$ is affine. Lemma~\ref{lem:ladder}
and $\sum_S|S|=kN/2$ give the gate bound.
Thresholding the negative of \eqref{eq:lookup} at zero computes $x_b$.
Bottom gates have address weights in $\{0,1\}$, input coefficients
of magnitude at most $2$, and thresholds of magnitude $O(n)$.
Finite differences of $\varphi$ are $O(n)$; combining repeated gates
preserves polynomial weight bounds.
\end{proof}

Appendix~\ref{sec:exact} saves gates by preserving only the score's sign.
For approximation, we truncate the staircases and bound the score error.

\section{Approximation from orthogonal odd errors}\label{sec:approx}
Throughout this section, all table and address bits are independent and
uniform. Recall $\varphi(u)=(-1)^u u$ for integer $u$. Balance means
$m=\lfloor N/k\rfloor$ or $km=N$, and logarithms are natural unless a
base is displayed.

\begin{theorem}[Uniform-input approximation]\label{thm:avgupper}
In the balanced case there are explicit sign-coherent depth-two circuits
with the following guarantees:
\begin{enumerate}
\item Error at most $\delta$ with
 $O(n^{3/2}\sqrt{\log(1/\delta)})$ gates, for $0<\delta\le1/2$.
\item Agreement at least $1/2+\varepsilon$ with
 $O(\varepsilon^2n^{3/2}\sqrt{\log(1/\varepsilon)})$ gates and
 $O(\varepsilon^2n^{5/2}\sqrt{\log(1/\varepsilon)})$ wires, for
 $N^{-1/2}\le\varepsilon\le1/8$.
\end{enumerate}
For $1/8<\varepsilon<1/2$, Part~1 applies with $\delta=1/2-\varepsilon$.
\end{theorem}

The following table symmetry will make the truncation errors orthogonal.
\begin{lemma}[Odd-error orthogonality]\label{lem:odd}
For a uniform independent table $x$, distinct $S,T\subseteq[k]$,
an odd function $g:\mathbb Z\to\mathbb R$, and any $h:\mathbb Z\to\mathbb R$,
$$\mathbb E_x[g(X_S)h(X_T)]=0.$$
The functions may depend on any fixed address input.
\end{lemma}
\begin{proof}
Choose $u\in\bits^k$ with $(S\mathbin\triangle T)\cdot u$ odd.
For $v_c=(-1)^{x_c}$, the signed permutation
$v'_c=(-1)^{T\cdot u}v_{c\oplus u}$ preserves the uniform distribution
of the table signs. It sends $X_S$ to $-X_S$ and leaves $X_T$ fixed.
Since $g$ is odd, the expectation equals its own negative and must vanish.
\end{proof}
The Walsh coefficients are not independent; the lemma uses their
pairwise sign symmetry.

\begin{proof}[Proof of Theorem~\ref{thm:avgupper}]
\emph{Partial lookup votes.}
For both parts, we approximate a vote over table bits containing $x_b$.
Choose $0\le r\le k$ and put $q=2^r$. Orthogonality gives
$$
P_*=2\sum_{S\subseteq[r]}X_S(-1)^{t_S}
 =2q\sum_{c:c_{[r]}=b_{[r]}}(-1)^{x_c}.
$$
The sum on the right contains $L=N/q$ table entries, selected by the
first $r$ coordinates of $b$. For $q=N$, the vote is exact. Otherwise
$L$ is even, and majority with ties outputting zero has agreement
$1/2+p(L)/2$, where
$$
p(L)=2^{-L}\binom L{L/2}\ge(2L)^{-1/2}.
$$
Conditioning on the selected bit and using symmetry, the excess
agreement is half the probability that the other $L-1$ signs sum to
$-1$; this probability is $p(L)$.
The inequality follows by induction from $p(2)=1/2$ and
$p(L+2)/p(L)=(L+1)/(L+2)\ge\sqrt{L/(L+2)}$.
In terms of scores, the vote is $[-P_*\ge2q]$ for $q<N$ and
$[-P_*\ge0]$ for $q=N$.

\emph{Symmetric truncation.}
Two reflected affine forms make the truncation error odd in $X_S$.
For nonempty $S\subseteq[r]$, set
$$D_S=t_S-\lfloor m|S|/2\rfloor,\qquad
\eta_S=(-1)^{\lfloor m|S|/2\rfloor}.$$
Since $X_S$ is even,
$$
\eta_S\bigl(\varphi(X_S+D_S)+\varphi(X_S-D_S)\bigr)
 =2X_S(-1)^{t_S}.
$$
For an integer $R\ge1$, truncate by setting
$f_R(u)=\varphi(u)[|u|\le R]$ and define
$$
P_R=2X_\varnothing+
\sum_{\varnothing\ne S\subseteq[r]}
\eta_S\bigl(f_R(X_S+D_S)+f_R(X_S-D_S)\bigr).
$$
For every integer $U$, the staircase identity
$$
f_R(U)=\sum_{\ell=-R}^{R+1}
\bigl(f_R(\ell)-f_R(\ell-1)\bigr)[U\ge\ell]
$$
uses at most $2R+2$ thresholds. Thus $P_R$ requires at most
$4(q-1)(R+1)$ bottom gates. The address weights are $+1$ on the selected
blocks for $X_S+D_S$, and $-1$ there for $X_S-D_S$; all other address
weights are zero. The gates are therefore sign-coherent.

\emph{The mean-square estimate.}
Let $e_S$ be the truncated $S$-summand minus its exact value, for
$\varnothing\ne S\subseteq[r]$. Both $\varphi$ and $f_R$ are odd.
For every fixed address input, replacing $X_S$ by $-X_S$ interchanges
and negates the reflected terms, so $e_S$ is odd in $X_S$.
Lemma~\ref{lem:odd} cancels the cross terms and gives
\begin{equation}\label{eq:odd-l2}
\mathbb E_{x,a}(P_R-P_*)^2=\sum_{\varnothing\ne S\subseteq[r]}\mathbb E_{x,a}e_S^2.
\end{equation}
We can now bound the squared errors individually.

For $U=X_S\pm D_S$, we have $|\mathbb EU|\le1/2$. Its centered part
sums $N$ independent table signs and $m|S|$ independent address signs
of magnitude $1/2$. For $V=U-\mathbb EU$ and $v_S=N+m|S|/4$,
Hoeffding's bound gives
$\Pr(|V|>t)\le2e^{-t^2/(2v_S)}$.
Since $|U|>R$ implies $|V|>R/2$ and $U^2\le2V^2+1/2$, integration of
this tail bound yields
$$
\mathbb E[U^2[|U|>R]]
 \le(R^2+8v_S+1)e^{-R^2/(8v_S)}
 \le16v_Se^{-R^2/(16v_S)}.
$$
The last step uses $v_S\ge2$ and $1+t\le e^t$.
Now $|\varphi(U)-f_R(U)|=|U|[|U|>R]$. Applying
$(u+v)^2\le2u^2+2v^2$ to the two reflected errors, then using
\eqref{eq:odd-l2} and $v=N+rm/4$, gives
$$
\mathbb E(P_R-P_*)^2\le64(q-1)v e^{-R^2/(16v)}.
$$

\emph{From scores to agreement.}
For $q<N$, the score $-P_*$ lies in $4q\mathbb Z$, at distance at
least $2q$ from the threshold $2q$. For $q=N$, the scores $\pm2N$
have the same distance from zero. Thus truncation changes the vote
only if $|P_R-P_*|\ge2q$, with probability at most
$$
\frac{16v}{q}e^{-R^2/(16v)}.
$$
Taking
$$R=\left\lceil4\sqrt{v\log\frac{16v}{q\delta}}\right\rceil$$
makes this probability at most $\delta$. This specifies a deterministic circuit. We may cap $R$ at $N+km$,
where truncation is exact on all attainable forms.

For Part~1, take $q=N$. The vote is exact, and balance gives
$v=O(N)$ and $R=O(\sqrt{N\log(1/\delta)})$.
For Part~2, let $q$ be the least power of two at least
$32\varepsilon^2N$ and set $\delta=\varepsilon$. Then
$q<64\varepsilon^2N$ and $q\le N$. If $q<N$, the vote's advantage is
at least $\sqrt{q/(8N)}\ge2\varepsilon$; if $q=N$, it is $1/2$.
Losing at most $\varepsilon$ through truncation leaves advantage at
least $\varepsilon$. Since
$\log(16v/(q\delta))=O(\log(1/\varepsilon))$, the stated gate bound
follows. Finally, a circuit with $G$ bottom gates has at most
$nG+G+n$ wires, including direct input wires to the output. In the
stated parameter range, the gate bound absorbs the additive terms and
gives the wire bound.
\end{proof}

\section{Restrictions of sign-coherent circuits}\label{sec:coherent}
Each restriction below leaves enough address inputs to read the whole
table, forcing many bottom gates to remain nonconstant. Sign coherence
limits the number of restrictions on which one gate can be nonconstant.
Comparing these counts gives the lower bound.

\begin{lemma}[Information preserved by a restriction]\label{lem:entropy}
Consider a fixed depth-two threshold circuit with $G$ bottom gates, input
table $x\in\bits^N$, and $h\ge1$ additional input bits $z$. Suppose the
function of $z$ computed for each $x$ determines $x$ uniquely. Let $D(x)$
be the number of bottom gates that are nonconstant as functions of $z$.
For uniform $x$,
$$
\mathbb E D(x)\ge
\frac{\max\{0,N-h-1-\log_2(G+1)\}}{h+\log_2(G+1)}.
$$
All weights may be arbitrary real numbers, and the output gate may read
both $x$ and $z$ directly.
\end{lemma}
\begin{proof}
Fixing $x$ shifts each bottom gate's threshold without changing its
weights on $z$. Moving this threshold gives at most $2^h-1$ nonconstant
functions on $\bits^h$. Specify the $D$ nonconstant gates by their
indices and one such function for each. Their output weights and the
weights on direct $z$-inputs are fixed. Constant bottom outputs and
direct table inputs only shift the output threshold, so at most $2^h+1$
output functions remain possible. Specifying which one determines $x$.
Using base-two entropy,
\begin{align*}
N=H(x)&\le H(D)+\mathbb E\log_2\left[
 \binom GD(2^h-1)^D(2^h+1)\right]\\
&\le\log_2(G+1)+\mathbb E D\bigl(\log_2(G+1)+h\bigr)+h+1.
\end{align*}
Rearranging and using $D\ge0$ proves the claim.
\end{proof}

\begin{theorem}[Exact sign-coherent lower bound]\label{thm:coherent}
For $N=2^k$, every exact sign-coherent circuit for $A_{k,m}$ with $G$
bottom gates satisfies
$$
G\ge\frac{D_0}{1-(1-1/m)^k}\ge\frac m kD_0,
\qquad
D_0=\frac{\max\{0,N-k-1-\log_2(G+1)\}}{k+\log_2(G+1)}.
$$
At floor or exact balance this is $\Omega(n^2/\log^3 n)$.
\end{theorem}
\begin{proof}
\emph{Each restriction reveals the table.}
Order the bits within each address block. For
$s\in\{0,\ldots,m\}^k$, let $a(s)$ have ones in exactly the first
$s_i$ positions of block $i$. Each $q\in\{0,\ldots,m-1\}^k$ defines
a cube of inputs $a(q+z)$, $z\in\bits^k$, with one live bit per block.
On this cube,
$$A_{k,m}(x,a(q+z))=x_{(q\bmod2)\oplus z}.$$
As $z$ varies, the output reads every table entry. Say that a gate
\emph{splits} a cube if it is nonconstant there. By
Lemma~\ref{lem:entropy} with $h=k$, every fixed cube has at least $D_0$
splitting gates in expectation over the table.

\emph{A gate splits few cubes.}
Fix a table and a bottom gate. In each block choose a direction
$\epsilon_i\in\{-1,1\}$ in which the gate's score is nondecreasing;
sign coherence guarantees such a direction. Let $F(s)$ denote its
score on $a(s)$, and let $\theta$ be its threshold. Unequal weight
magnitudes and zero weights are allowed. On the cube indexed by $q$,
the score attains its minimum and maximum at
$$q^-=q+(\mathbf1-\epsilon)/2,\qquad
q^+=q+(\mathbf1+\epsilon)/2.$$
The gate splits that cube exactly when
$F(q^-)<\theta\le F(q^+)$.

Partition the cube indices into lines with successive indices differing
by $\epsilon$. The maximum vertex of one cube is the minimum vertex
of the next. Hence the score intervals $(F(q^-),F(q^+)]$ along a line
are disjoint, and a gate splits at most one cube on each line.
A line starts exactly where $q-\epsilon$ lies outside the grid.
There are $(m-1)^k$ indices with a predecessor, so there are
$m^k-(m-1)^k$ lines.

\emph{Combining the counts.}
Summing over the $m^k$ cubes gives the expected number of incidences
between gates and cubes they split. Each gate contributes at most
$m^k-(m-1)^k$ for every table, so
$$m^kD_0\le G\bigl(m^k-(m-1)^k\bigr).$$
Dividing gives the first inequality, and
$1-(1-1/m)^k\le k/m$ gives the second.
For the asymptotic claim, $G\ge N^2$ already suffices. Otherwise,
integrality gives $G+1\le N^2$, whence
$D_0\ge(N-3k-1)/(3k)$ for large $k$. At balance,
$m=\Theta(N/k)$ and $n=\Theta(N)$, which proves the stated bound.
\end{proof}

\paragraph{Why exactness matters.}
These restrictions range over block counts, not uniform raw addresses,
so a uniform-input error bound need not survive them. This allows the
sign-coherent approximation circuits of Section~\ref{sec:approx}.
An exact construction with exponent below $2$ by a fixed amount must
use mixed address-weight signs within a block of some gate.

\section{Further questions}
For ordinary Andreev lookup, the unrestricted exact gate bounds still
have exponents $3/2$ and $2$. Can mixed signs within address blocks
yield $n^{2-\eta}$ gates for fixed $\eta>0$? Can the exact wire bound
of $O(n^3)$ improve toward the known lower exponent $5/2$? Gate upper
bounds imply wire upper bounds with a factor of $O(n)$; wire lower
bounds require information about which inputs each gate uses.

At fixed error in $(0,1/2)$, the gate exponent is determined up to
logarithmic factors, but the error dependence remains open. Our bound
becomes quadratic when $\log(1/\delta)=\Theta(n)$. For small advantage,
the upper bound's $\varepsilon^2$ dependence differs from the lower
bound's $\varepsilon^3$, where that lower bound applies. Closing either
gap would clarify the transition from weak prediction to exact computation.
Do structured or biased table distributions retain enough of the sign
symmetry in Lemma~\ref{lem:odd} for the same approximation argument?

\paragraph{Verification and assistance.}
The research and exposition were developed with AI assistance.
Finite checks accompany the source: \texttt{verify\_results.py} checks
the refined construction; \texttt{approximation\_checks.py} checks
truncation and orthogonality; and \texttt{restriction\_checks.py} checks
the restrictions. They supplement the proofs. An earlier quadratic
construction is formalized in Lean in
\href{https://github.com/cslib-community/cslib-frontier/pull/2813}{cslib-frontier};
no formalization is claimed for the other results here.

\appendix
\section{The sharper exact construction}\label{sec:exact}
We improve the constant in Theorem~\ref{thm:simple} using one staircase
per nonempty character, with thresholds only at odd levels. Its
quadratic corrections combine into one global square and $k$ block-count
squares. Further corrections handle the zero and singleton addresses.
We verify the score's sign, then count its gates. All auxiliary notation
is local to this appendix.

\begin{theorem}\label{thm:exact}
Put $\alpha=\lfloor(m-1)/2\rfloor$. There is an exact sign-coherent
depth-two threshold circuit for $A_{k,m}$ with at most
\begin{equation}\label{eq:bestcount}
\frac{N(N-1)}2+\frac{kN(m+\alpha)}2
 +2(N-1)+2\alpha k+3km+\left\lceil\frac{km}{2}\right\rceil
\end{equation}
bottom gates and one output gate. Its weights are integers of magnitude
polynomial in $n$. At exact balance its total size is
$$
\frac5{16}n^2-\frac n4\log_2(n/2)+O(n).
$$
For $m=\lfloor N/k\rfloor$, its size is $(5/16+o(1))n^2$.
\end{theorem}

\begin{proof}
\emph{Affine forms and the staircase.}
Write $s_i=|a_i|$ and $\sigma=\sum_i s_i$, so
$t_S=\sum_{i\in S}s_i$. Let $e_i$ be the address with only coordinate
$i$ equal to one, and let $x_0$ denote the entry at the all-zero address.
Set
$$
z=\sum_{c\ne0}x_c,\quad E=\sum_i x_{e_i},\quad
w=2z+2\alpha E+\sigma,\quad R=2(N-1)+2\alpha k+km,\quad C=N(R+4).
$$
Thus $0\le w\le R$. For each nonempty $S\subseteq[k]$, we use a sum over the
$N/2$ table positions with $S\cdot c$ odd, augmented by a correction
for singleton addresses. Define
$$
W_S=\sum_{c:S\cdot c\text{ odd}}x_c+\alpha\sum_{i\in S}x_{e_i},\quad
\kappa_S=-4(2\alpha+1)(k-2|S|),
$$
and put $r_S(u)=-8\lceil u/2\rceil^2+\kappa_S\lceil u/2\rceil$.
The forms $W_S,t_S,w,s_i$ are all affine in the input bits.
For integers $W,t$, the identity underlying the staircase is
$$
4W(-1)^t=4W-8\lceil(2W+t)/2\rceil^2+8W^2+8Wt+8\lceil t/2\rceil^2.
$$
The ceiling makes $r_S$ change only at odd integer levels, saving
thresholds. The shifts $\kappa_S$ correct singleton addresses using
those same thresholds.

\emph{The score and its sign.}
Consider
\begin{align}\label{eq:bestscore}
P={}&2N+\sum_{S\ne\varnothing}r_S(2W_S+t_S)
 +\frac N2\left(w^2+\sum_i s_i^2\right)\notag\\
 &+4Nz-2N(\alpha^2+1)E-2\alpha N\sigma
 +2N(1+\alpha)\sum_i x_{e_i}s_i-Cx_0\ind{b=0}.
\end{align}
The squares collect the quadratic corrections; the remaining terms
adjust the exceptional addresses. Expanding the identity above uses
three sums, all over nonempty $S$:
\begin{align*}
\sum_S4W_S(-1)^{t_S}
 &=2N\ind{b=0}z-2N\ind{b\ne0}x_b
   +2\alpha N\sum_i x_{e_i}\bigl(\ind{b=0}-\ind{b=e_i}\bigr),\\
8\sum_S(W_S^2+W_St_S)
 &=\frac N2(w^2-\sigma^2)+2Nz+(4\alpha+2\alpha^2)NE
   +2N(1+\alpha)\sum_i x_{e_i}s_i,\\
8\sum_S\lceil t_S/2\rceil^2
 &=\frac N2\left(\sigma^2+\sum_i s_i^2\right)
   +N(\sigma+1)\ind{b\ne0}+N\sum_i s_i\ind{b=e_i}.
\end{align*}
The first identity follows from character orthogonality. For the
quadratic terms, distinct nonzero addresses $c,d$ satisfy
$$\sum_S[S\cdot c\text{ odd}]=N/2,\qquad
\sum_S[S\cdot c\text{ odd}][S\cdot d\text{ odd}]=N/4.$$
When $c=d$, the second count is $N/2$. These counts give the expansion
of $W_S^2$ and $W_St_S$. For the last sum, use
$\lceil t/2\rceil=(t+[t\text{ odd}])/2$ and the same orthogonality.
The shifts can be summed using $k-2|S|=\sum_i(-1)^{[i\in S]}$.
Their parity-dependent contributions are
$$
\sum_S\kappa_S\bigl(\lceil t_S/2\rceil-t_S/2\bigr)
 =N(2\alpha+1)\sum_i\ind{b=e_i},
$$
and their linear contributions are
$$\sum_S\kappa_SW_S=2N(2\alpha+1)(1+\alpha)E,\qquad
\frac12\sum_S\kappa_St_S=N(2\alpha+1)\sigma.$$
Substituting these identities in \eqref{eq:bestscore} gives
\begin{equation}\label{eq:bestclosed}
P=\begin{cases}
N(w+2)-Cx_0,&b=0,\\
2N(1+\alpha)(1-x_{e_i})-Ns_i,&b=e_i,\\
N(1-2x_b),&\text{otherwise}.
\end{cases}
\end{equation}
For $b=0$, use $0\le w\le R$ and $C=N(R+4)$. For $b=e_i$, the odd
count $s_i$ lies between $1$ and $2\alpha+1$. In every case, $P\ge N$
when $x_b=0$ and $P\le-N$ when $x_b=1$, so the output is
$\ind{-P\ge0}$.

\emph{Realizing and counting the gates.}
The staircase $r_S$ changes only at odd integers. For
$r_S(2W_S+t_S)$, use the gates
$$
\ind{2W_S+t_S\ge\theta},\qquad
1\le\theta\le N+2\alpha|S|+|S|m,\quad \theta\text{ odd},
$$
with output coefficient $-8\theta+\kappa_S$. This also explains why
the shifts $\kappa_S$ add no gates. By Lemma~\ref{lem:ladder}, $w^2$
and the block squares require $R+km$ gates in total.
For $B=n+2$, each product $x_{e_i}s_i$ is a sum of $m$ gates
$\ind{s_i+Bx_{e_i}\ge B+\ell}$, $1\le\ell\le m$.

To implement the correction at the zero address, use
$$
x_0\ind{b=0}=x_0-x_0\sigma+
 \frac4N\sum_{S\ne\varnothing}x_0\lfloor t_S/2\rfloor.
$$
This follows from $\sum_S\ind{t_S\text{ odd}}=(N/2)\ind{b\ne0}$
and $\sum_S t_S=N\sigma/2$. The product $x_0\sigma$ uses $km$ gates
$\ind{\sigma+Bx_0\ge B+\ell}$, $1\le\ell\le km$.
The remaining products use gates
$\ind{t_S+Bx_0\ge B+2j}$ for
$1\le j\le\lfloor |S|m/2\rfloor$.
The even-level $\sigma$ gates coincide with those for $S=[k]$, saving
$\lfloor km/2\rfloor$ gates. Every other term in $P$ is affine and
can be supplied directly to the output. All address weights in these
bottom gates are nonnegative.

Since $\sum_{S\ne\varnothing}|S|=kN/2$, the character staircases
and the $x_0\lfloor t_S/2\rfloor$ terms together cost at most
$$
\sum_{S\ne\varnothing}\left(
 \left\lceil\frac{N+2\alpha|S|+|S|m}{2}\right\rceil
 +\left\lfloor\frac{|S|m}{2}\right\rfloor\right)
 =\frac{N(N-1)}2+\frac{\alpha kN}2+\frac{kmN}2.
$$
Adding the remaining gates and subtracting shared gates gives
\eqref{eq:bestcount}. Zero coefficients or further coincidences only
reduce the count. At exact balance, $m$ is even and
$\alpha=m/2-1$; substitution gives the stated size. At floor balance,
rounding changes the leading expression by $O(Nk)=o(n^2)$.
Bottom coefficients and thresholds are $O(n)$, and output coefficients
in the expansion are $O(n^3)$. All coefficients are integers, since
$N$ is even and $C$ is divisible by $N$.
\end{proof}

\section{A full-support equal-magnitude obstruction}\label{sec:fullsign}
Here every nonzero address-weight vector must have full support and
equal magnitudes. Each of exponentially many restrictions requires a
gate depending on all live variables, by the classical parity-order
obstruction~\cite[Section~3.1]{MP}. A gate can meet this requirement
on at most one restriction.

\begin{lemma}[The missing-variable parity obstruction]\label{lem:paritydual}
Let $d\ge2$ and $f(z)=\prod_{i=1}^d z_i$ on $\{-1,1\}^d$.
If a score consisting of an affine term and a linear combination of
predicates has the same strict sign as $f$ at every input, then some
predicate depends on every variable.
\end{lemma}
\begin{proof}
A predicate missing coordinate $i$ has zero correlation with $f$:
summing over $z_i$ cancels. Affine terms also have zero correlation
when $d\ge2$. If every predicate omitted a variable, the score $F$
would satisfy $\mathbb E[fF]=0$, contradicting strict sign agreement,
which gives $\mathbb E[fF]=\mathbb E|F|>0$.
\end{proof}

\begin{proposition}\label{thm:fullsign}
Let $k\ge2$, $m\ge1$, and $d=m\lfloor k/2\rfloor\ge2$. Suppose every
bottom gate's vector of address weights is either zero or $\rho\sigma$,
where $\rho>0$ and $\sigma\in\{-1,1\}^{km}$ may depend on the gate.
An exact circuit for $A_{k,m}$ then has at least $2^d$ bottom gates.
Table weights, thresholds, output weights, and direct input wires to
the output are unrestricted; all weights may be real.
At balance the bound is $2^{\Omega(n)}$.
\end{proposition}
\begin{proof}
Fix the table so that its entry at address $c$ is the parity of the
first $r=\lfloor k/2\rfloor$ coordinates of $c$. If $k$ is odd, fix
one of the other address blocks to zero. There remain $d=rm$ selected
bits and $d$ ignored bits. Pair them and write their sign coordinates
$2a-1$ as $z_i,w_i$. Up to a constant sign, the target is
$f(z)=\prod_i z_i$.
Write the circuit's output as $[F\ge0]$. On inputs with output bit zero,
$F<0$. Since the domain is finite, adding a sufficiently small positive
constant makes the score nonzero everywhere without changing the
output. Orient this strict score to match $f$.

For each $t\in\{-1,1\}^d$, impose $w_i=t_i z_i$.
These bit or complemented-bit identifications keep direct output
terms affine.
A gate with nonzero address weights and paired signs $\sigma_i,\tau_i$
then has coefficients proportional to $\sigma_i+t_i\tau_i$.
It can depend on every $z_i$ only if $t_i=\sigma_i\tau_i$ for all $i$,
which specifies a unique restriction. A gate with zero address weights
is constant after the table is fixed.
Lemma~\ref{lem:paritydual} requires a gate depending on all live variables
for each of the $2^d$ restrictions. Each original gate can serve at most
one of them, proving the bound.
\end{proof}

Both qualifications are necessary. Without full support,
Theorem~\ref{thm:simple} gives quadratic size with address weights in
$\{0,1\}$. Without equal magnitudes, recall $X_S,t_S,\varphi$ from
Section~\ref{sec:identities} and put
$$L_S=\sum_{i\in S,j}a_{i,j}+2\sum_{i\notin S,j}a_{i,j}.$$
The form $L_S$ has full address support and the parity of $t_S$.
Thus $X_S(-1)^{t_S}=\varphi(X_S+L_S)-\varphi(L_S)$, and
Lemma~\ref{lem:ladder} gives $O(n^2)$ gates because each range has
length $O(n)$. The exponential bound does not extend to full-support
address weights of general bounded integer magnitudes.

\begin{thebibliography}{9}\small
\bibitem{OD} R.~O'Donnell. \emph{Analysis of Boolean Functions}.
Cambridge University Press, 2014.
\href{https://www.cs.cmu.edu/~odonnell/papers/Analysis-of-Boolean-Functions-by-Ryan-ODonnell.pdf}{Author-hosted edition}.
\bibitem{CTW} L.~Chen, A.~Tal, and Y.~Wang.
Super-quadratic Lower Bounds for Depth-2 Linear Threshold Circuits.
\href{https://eccc.weizmann.ac.il/report/2026/039/}{ECCC TR26-039}, 2026.
\bibitem{KW} D.~M. Kane and R.~Williams.
Super-Linear Gate and Super-Quadratic Wire Lower Bounds for Depth-Two and
Depth-Three Threshold Circuits. STOC, 2016.
\href{https://arxiv.org/abs/1511.07860}{arXiv version: 1511.07860}.
\bibitem{HMPST} A.~Hajnal, W.~Maass, P.~Pudl\'ak, M.~Szegedy, and G.~Tur\'an.
Threshold circuits of bounded depth.
\emph{Journal of Computer and System Sciences} 46(2):129--154, 1993.
\bibitem{HP} K.~A. Hansen and V.~V. Podolskii.
Polynomial threshold functions and Boolean threshold circuits.
\emph{Information and Computation} 240:56--73, 2015.
References to theorem and lemma numbers use the
\href{https://eccc.weizmann.ac.il/report/2013/021/download/}{ECCC version, TR13-021}.
\bibitem{MP} M.~L. Minsky and S.~A. Papert.
\emph{Perceptrons: An Introduction to Computational Geometry}.
MIT Press, 1969; expanded edition, 1988. The parity-order argument is in
Section~3.1.
\bibitem{IPS} R.~Impagliazzo, R.~Paturi, and M.~Saks.
Size--depth tradeoffs for threshold circuits.
\emph{SIAM Journal on Computing} 26(3):693--707, 1997.
\bibitem{CSS} R.~Chen, R.~Santhanam, and S.~Srinivasan.
Average-Case Lower Bounds and Satisfiability Algorithms for Small Threshold
Circuits. \emph{Theory of Computing} 14(9):1--55, 2018.
\end{thebibliography}
\end{document}
